% =============================================================================
% Chapter 9 — The File I/O Controller ($B9A0–$B9EF)
% =============================================================================
\chapter{The File I/O Controller}

\epigraph{``A computer without persistent storage is just a very expensive
calculator.''}{\textit{--- Unknown, circa 1975}}

\section*{Introduction}

The File I/O Controller --- FIO --- is a coprocessor that bridges the 6502's
address space to the host filesystem.  It handles file save/load, directory
enumeration, SID file playback, music sequencing via MML, MIDI file playback,
soundfont loading, and virtual disk operations.

Commands are issued by writing to the command register at \$B9A0.  Parameters
are set in surrounding registers \emph{before} the command write.  After
execution, the status register at \$B9A1 reports success or failure, and the
error code register at \$B9A2 provides detail.

Files are stored in the save directory (\texttt{\textasciitilde/e6502-programs}
by default).  The FIO automatically appends the appropriate file extension
(\texttt{.bas}, \texttt{.sid}, \texttt{.gfx}, etc.) based on the command.


% ═══════════════════════════════════════════════════════════════════════════════
\section{Register Map}
\label{sec:fio-registers}
% ═══════════════════════════════════════════════════════════════════════════════

\begin{center}
\small
\begin{tabular}{@{}llcll@{}}
\toprule
\textbf{Address} & \textbf{Name} & \textbf{R/W} & \textbf{Default} & \textbf{Description} \\
\midrule
\$B9A0 & FioCmd      & W   & ---  & Command register (write triggers execution) \\
\$B9A1 & FioStatus   & R   & \$00 & Status code \\
\$B9A2 & FioErrCode  & R   & \$00 & Error code \\
\$B9A3 & FioNameLen  & R/W & 0    & Filename length (1--63) \\
\$B9A4 & FioSrcL     & R/W & 0    & Source/destination RAM address (low byte) \\
\$B9A5 & FioSrcH     & R/W & 0    & Source/destination RAM address (high byte) \\
\$B9A6 & FioEndL     & R/W & 0    & End address for SAVE (low byte) \\
\$B9A7 & FioEndH     & R/W & 0    & End address for SAVE (high byte) \\
\$B9A8 & FioSizeL    & R   & 0    & Loaded data size (low byte) \\
\$B9A9 & FioSizeH    & R   & 0    & Loaded data size (high byte) \\
\$B9AA & FioGSpace   & R/W & 0    & Graphics space selector \\
\$B9AB & FioGAddrL   & R/W & 0    & Graphics offset (low byte) \\
\$B9AC & FioGAddrH   & R/W & 0    & Graphics offset (high byte) \\
\$B9AD & FioGLenL    & R/W & 0    & Graphics length (low byte) \\
\$B9AE & FioGLenH    & R/W & 0    & Graphics length (high byte) \\
\$B9AF & FioDirType  & R   & 0    & Directory entry file type \\
\$B9B0--\$B9EF & FioName & R/W & 0 & Filename buffer (64 bytes ASCII) \\
\bottomrule
\end{tabular}
\end{center}


% ═══════════════════════════════════════════════════════════════════════════════
\section{Core Registers}
\label{sec:fio-core}
% ═══════════════════════════════════════════════════════════════════════════════

% --- $B9A0 FioCmd ----------------------------------------------------------

\mementry{B9A0}{FioCmd}{W}{---}

Command register.  Writing a command byte here immediately executes the
corresponding operation.  All parameters (filename, addresses, lengths) must be
set in the other registers \emph{before} writing to FioCmd.

% --- $B9A1 FioStatus -------------------------------------------------------

\mementry{B9A1}{FioStatus}{R}{\$00}

Status of the last command:

\begin{center}
\small
\begin{tabular}{@{}cl@{}}
\toprule
\textbf{Value} & \textbf{Meaning} \\
\midrule
\$00 & Idle --- no command has been issued \\
\$02 & OK --- last command completed successfully \\
\$03 & Error --- last command failed (see FioErrCode) \\
\bottomrule
\end{tabular}
\end{center}

% --- $B9A2 FioErrCode ------------------------------------------------------

\mementry{B9A2}{FioErrCode}{R}{\$00}

Detailed error code when FioStatus is \$03:

\begin{center}
\small
\begin{tabular}{@{}cl@{}}
\toprule
\textbf{Value} & \textbf{Meaning} \\
\midrule
0 & No error \\
1 & File not found \\
2 & I/O error \\
3 & End of directory \\
4 & Disk full \\
5 & Device not mounted \\
\bottomrule
\end{tabular}
\end{center}

% --- $B9A3 FioNameLen -------------------------------------------------------

\mementry{B9A3}{FioNameLen}{R/W}{0}

Length of the filename stored in the FioName buffer.  Valid range: 1--63.
Some commands use this value; others read the buffer directly until they
encounter a null byte or reach 64 characters.

% --- $B9A4-$B9A5 FioSrc ---------------------------------------------------

\mementry{B9A4}{FioSrcL}{R/W}{0}

Low byte of the 16-bit source or destination address in CPU RAM.  For
\textbf{Save}, this is the start of the region to save.  For \textbf{Load},
this is where loaded data will be written.  For music commands, the low byte
may carry a parameter value (e.g., instrument index, volume level).

\mementry{B9A5}{FioSrcH}{R/W}{0}

High byte of the source/destination address.

% --- $B9A6-$B9A7 FioEnd ---------------------------------------------------

\mementry{B9A6}{FioEndL}{R/W}{0}

Low byte of the end address.  For \textbf{Save}, the region
FioSrc--FioEnd (inclusive) is written to disk.  For \textbf{MusicSeq}, FioEnd
minus FioSrc gives the length of the MML string in RAM.

\mementry{B9A7}{FioEndH}{R/W}{0}

High byte of the end address.

% --- $B9A8-$B9A9 FioSize --------------------------------------------------

\mementry{B9A8}{FioSizeL}{R}{0}

Low byte of the loaded data size.  Set by the host after a successful
\textbf{Load} command.

\mementry{B9A9}{FioSizeH}{R}{0}

High byte of the loaded data size.

% --- $B9AA FioGSpace -------------------------------------------------------

\mementry{B9AA}{FioGSpace}{R/W}{0}

Graphics space selector for \textbf{GSave}/\textbf{GLoad} commands:

\begin{center}
\small
\begin{tabular}{@{}cl@{}}
\toprule
\textbf{Value} & \textbf{Space} \\
\midrule
0 & Screen (Character RAM) \\
1 & Color (Color RAM) \\
2 & Graphics (bitmap pixel data) \\
3 & Sprite (sprite definition data) \\
\bottomrule
\end{tabular}
\end{center}

% --- $B9AB-$B9AC FioGAddr --------------------------------------------------

\mementry{B9AB}{FioGAddrL}{R/W}{0}

Low byte of the offset within the selected graphics space.

\mementry{B9AC}{FioGAddrH}{R/W}{0}

High byte of the graphics offset.

% --- $B9AD-$B9AE FioGLen --------------------------------------------------

\mementry{B9AD}{FioGLenL}{R/W}{0}

Low byte of the graphics data length.

\mementry{B9AE}{FioGLenH}{R/W}{0}

High byte of the graphics data length.

% --- $B9AF FioDirType ------------------------------------------------------

\mementry{B9AF}{FioDirType}{R}{0}

File type of the current directory entry, populated by \textbf{DirRead}:

\begin{center}
\small
\begin{tabular}{@{}cl@{}}
\toprule
\textbf{Value} & \textbf{Type} \\
\midrule
0 & \texttt{.bas} --- BASIC program \\
1 & \texttt{.sid} --- SID music file \\
2 & \texttt{.bin} --- Binary executable \\
3 & \texttt{.mid} --- MIDI file \\
5 & Directory \\
\bottomrule
\end{tabular}
\end{center}

% --- $B9B0-$B9EF FioName --------------------------------------------------

\memrange{B9B0}{B9EF}{FioName}{R/W}{0}

64-byte filename buffer.  Store the filename as ASCII characters starting at
\$B9B0 before issuing commands that require a filename.  Null-terminate the
string or set FioNameLen to the correct length.

\begin{lstlisting}
10 REM store filename "MYFILE" into FioName buffer
20 F$ = "MYFILE" + CHR$(0)
30 FOR I=0 TO LEN(F$)-1
40   POKE $B9B0+I, ASC(MID$(F$,I+1,1))
50 NEXT I
\end{lstlisting}


% ═══════════════════════════════════════════════════════════════════════════════
\section{Command Reference}
\label{sec:fio-commands}
% ═══════════════════════════════════════════════════════════════════════════════

Commands are executed by writing the command byte to \addr{B9A0}.  Set all
required parameters \emph{before} writing the command byte.

% ─────────────────────────────────────────────────────────────────────────────
\subsection{File Operations}
% ─────────────────────────────────────────────────────────────────────────────

\cmdentry{01}{Save}{Save BASIC program to disk}

Saves the region of CPU RAM from FioSrc to FioEnd (inclusive) as a
\texttt{.bas} file.  The filename is read from the FioName buffer.

\textbf{Parameters:}
\begin{itemize}
  \item FioSrc (\$B9A4--\$B9A5): start address of data in RAM
  \item FioEnd (\$B9A6--\$B9A7): end address of data in RAM
  \item FioName (\$B9B0+): filename (without extension)
\end{itemize}

\textbf{Result:} FioStatus $=$ \$02 on success, \$03 on error.

\medskip

\cmdentry{02}{Load}{Load BASIC program from disk}

Loads a \texttt{.bas} file into CPU RAM starting at FioSrc.  After a
successful load, FioSize (\$B9A8--\$B9A9) contains the number of bytes
loaded.

\textbf{Parameters:}
\begin{itemize}
  \item FioSrc (\$B9A4--\$B9A5): destination address in RAM
  \item FioName (\$B9B0+): filename (without extension)
\end{itemize}

\textbf{Result:} FioStatus $=$ \$02, FioSize updated; or \$03 with
ErrCode $=$ 1 (not found) or 2 (I/O error).

\medskip

\cmdentry{03}{DirOpen}{Open directory for enumeration}

Opens the current save directory for sequential reading with DirRead.
Resets the internal directory index to the first entry.

\textbf{Parameters:} None.

\textbf{Result:} FioStatus $=$ \$02.

\medskip

\cmdentry{04}{DirRead}{Read next directory entry}

Reads the next file in the directory.  The filename is placed in the FioName
buffer and the file type in FioDirType (\$B9AF).  When all entries have been
read, FioStatus is set to \$03 with ErrCode $=$ 3 (end of directory).

\textbf{Parameters:} None (directory must be opened with DirOpen first).

\textbf{Result:} FioName and FioDirType populated, or ErrCode $=$ 3 at end.

\medskip

\cmdentry{05}{Delete}{Delete a file}

Deletes the file named in the FioName buffer.

\textbf{Parameters:}
\begin{itemize}
  \item FioName (\$B9B0+): filename (without extension)
\end{itemize}

\textbf{Result:} FioStatus $=$ \$02 on success, \$03 with ErrCode $=$ 1 if
not found.

\medskip

\cmdentry{06}{GSave}{Save VGC memory to disk}

Saves a region of VGC memory to a \texttt{.gfx} file.  The space selector
determines which VGC buffer is read.

\textbf{Parameters:}
\begin{itemize}
  \item FioGSpace (\$B9AA): graphics space (0--3)
  \item FioGAddr (\$B9AB--\$B9AC): offset within the space
  \item FioGLen (\$B9AD--\$B9AE): number of bytes to save
  \item FioName (\$B9B0+): filename
\end{itemize}

\textbf{Result:} FioStatus $=$ \$02 on success.

\medskip

\cmdentry{07}{GLoad}{Load VGC memory from disk}

Loads a \texttt{.gfx} file into the selected VGC memory space.

\textbf{Parameters:}
\begin{itemize}
  \item FioGSpace (\$B9AA): target graphics space (0--3)
  \item FioGAddr (\$B9AB--\$B9AC): offset within the space
  \item FioGLen (\$B9AD--\$B9AE): number of bytes to load
  \item FioName (\$B9B0+): filename
\end{itemize}

\textbf{Result:} FioStatus $=$ \$02 on success.


% ─────────────────────────────────────────────────────────────────────────────
\subsection{SID Playback}
% ─────────────────────────────────────────────────────────────────────────────

\cmdentry{08}{SidPlay}{Play a .sid file}

Loads and plays a \texttt{.sid} file through the SID chip emulation.  The
SID player injects an IRQ trampoline into CPU RAM that calls the file's
init and play routines at 60\,Hz.

\textbf{Parameters:}
\begin{itemize}
  \item FioName (\$B9B0+): filename (without \texttt{.sid} extension)
\end{itemize}

\textbf{Result:} FioStatus $=$ \$02 if playback started.

\medskip

\cmdentry{09}{SidStop}{Stop SID playback}

Stops the currently playing SID file and removes the IRQ trampoline.

\textbf{Parameters:} None.

\textbf{Result:} FioStatus $=$ \$02.


% ─────────────────────────────────────────────────────────────────────────────
\subsection{Music Engine}
% ─────────────────────────────────────────────────────────────────────────────

The Music Engine provides a 14-voice sequencer (6 SID + 8 WTS) driven by
MML (Music Macro Language) strings.  These commands control instrument
selection, sound effects, volume, and MML sequence playback.

\cmdentry{0A}{Instrument}{Select SID instrument}

Selects a SID instrument preset for subsequent sound and music commands.

\textbf{Parameters:}
\begin{itemize}
  \item FioSrcL (\$B9A4): instrument index (0--15)
\end{itemize}

\textbf{Result:} FioStatus $=$ \$02.

\medskip

\cmdentry{0B}{Sound}{Trigger a sound effect}

Triggers a one-shot sound effect on the SID, using the currently selected
instrument.  The effect uses voice stealing if all voices are busy.

\textbf{Parameters:}
\begin{itemize}
  \item FioSrcL (\$B9A4): MIDI note number (0--127)
\end{itemize}

\textbf{Result:} FioStatus $=$ \$02.

\medskip

\cmdentry{0C}{Volume}{Set music volume}

Sets the master volume for the music engine.

\textbf{Parameters:}
\begin{itemize}
  \item FioSrcL (\$B9A4): volume level (0--255)
\end{itemize}

\textbf{Result:} FioStatus $=$ \$02.

\medskip

\cmdentry{0D}{MusicSeq}{Load MML music sequence}

Loads an MML string from CPU RAM into the music engine for playback.  The
MML string is read from the address range FioSrc to FioEnd.

\textbf{Parameters:}
\begin{itemize}
  \item FioSrc (\$B9A4--\$B9A5): start address of MML string in RAM
  \item FioEnd (\$B9A6--\$B9A7): end address of MML string in RAM
\end{itemize}

\textbf{Result:} FioStatus $=$ \$02 if the MML was parsed successfully.

\begin{notebox}
In practice, the BASIC \texttt{MUSIC} statement is the easiest way to load
and play MML sequences.  Direct FIO access is useful when MML data is
generated programmatically or stored in expansion memory.
\end{notebox}

\medskip

\cmdentry{0E}{MusicPlay}{Start music playback}

Begins playing the MML sequence loaded by MusicSeq.

\textbf{Parameters:} None.

\textbf{Result:} FioStatus $=$ \$02.

\medskip

\cmdentry{0F}{MusicStop}{Stop music playback}

Stops the currently playing music sequence and silences all music voices.

\textbf{Parameters:} None.

\textbf{Result:} FioStatus $=$ \$02.

\medskip

\cmdentry{10}{MusicTempo}{Set music tempo}

Sets the tempo in beats per minute.

\textbf{Parameters:}
\begin{itemize}
  \item FioSrc (\$B9A4--\$B9A5): tempo in BPM (16-bit value)
\end{itemize}

\textbf{Result:} FioStatus $=$ \$02.

\medskip

\cmdentry{11}{MusicLoop}{Set loop mode}

Controls whether the music sequence loops or plays once.

\textbf{Parameters:}
\begin{itemize}
  \item FioSrcL (\$B9A4): 0 $=$ play once, 1 $=$ loop
\end{itemize}

\textbf{Result:} FioStatus $=$ \$02.

\medskip

\cmdentry{12}{MusicPri}{Set music priority}

Sets the priority level for the music engine, affecting voice-stealing
behaviour when sound effects compete with background music.

\textbf{Parameters:}
\begin{itemize}
  \item FioSrcL (\$B9A4): priority level
\end{itemize}

\textbf{Result:} FioStatus $=$ \$02.


% ─────────────────────────────────────────────────────────────────────────────
\subsection{MIDI Playback}
% ─────────────────────────────────────────────────────────────────────────────

\cmdentry{13}{MidPlay}{Play a MIDI file}

Loads and plays a standard \texttt{.mid} file through the Music Engine's
14-voice system.  Three routing modes are available: Auto (WTS preferred
with SID overflow), Manual (per-channel table), and SID-only (legacy
6-voice).  GM program numbers are automatically mapped to WTS instrument
indices or SID instrument buckets.

\textbf{Parameters:}
\begin{itemize}
  \item FioName (\$B9B0+): filename (without \texttt{.mid} extension)
\end{itemize}

\textbf{Result:} FioStatus $=$ \$02 if playback started.

\medskip

\cmdentry{14}{MidStop}{Stop MIDI playback}

Stops the currently playing MIDI file and releases all voices.

\textbf{Parameters:} None.

\textbf{Result:} FioStatus $=$ \$02.


% ─────────────────────────────────────────────────────────────────────────────
\subsection{Soundfont}
% ─────────────────────────────────────────────────────────────────────────────

\cmdentry{15}{SfLoad}{Load SF2 soundfont for WTS}

Loads an SF2 soundfont file from the host filesystem into the Wavetable
Synthesizer.  Sample data is kept in host memory.  The WTS status register
at \addr{A183} can be polled to detect when loading is complete.

\textbf{Parameters:}
\begin{itemize}
  \item FioName (\$B9B0+): soundfont filename (with \texttt{.sf2} extension)
\end{itemize}

\textbf{Result:} FioStatus $=$ \$02 if loading was initiated.  Poll
\texttt{PEEK(\$A183)} for ready status.


% ─────────────────────────────────────────────────────────────────────────────
\subsection{Disk Operations}
% ─────────────────────────────────────────────────────────────────────────────

These commands operate on the virtual storage device system.  They require a
mounted storage device; issuing them without one produces ErrCode $=$ 5
(not mounted).

\cmdentry{20}{Cd}{Change directory}

Changes the current working directory on the active storage device.

\textbf{Parameters:}
\begin{itemize}
  \item FioName (\$B9B0+): directory name or path
\end{itemize}

\textbf{Result:} FioStatus $=$ \$02 on success.

\medskip

\cmdentry{21}{Mkdir}{Create directory}

Creates a new directory.

\textbf{Parameters:}
\begin{itemize}
  \item FioName (\$B9B0+): directory name
\end{itemize}

\textbf{Result:} FioStatus $=$ \$02 on success.

\medskip

\cmdentry{22}{Rmdir}{Remove directory}

Removes an empty directory.

\textbf{Parameters:}
\begin{itemize}
  \item FioName (\$B9B0+): directory name
\end{itemize}

\textbf{Result:} FioStatus $=$ \$02 on success.

\medskip

\cmdentry{23}{Format}{Format disk}

Formats the active storage device, erasing all data.

\textbf{Parameters:} None.

\textbf{Result:} FioStatus $=$ \$02 on success.

\begin{warningbox}
Format erases all data on the device.  There is no confirmation prompt
and no undo.
\end{warningbox}

\medskip

\cmdentry{24}{Mount}{Mount disk}

Mounts a storage device, making it available for file operations.

\textbf{Parameters:}
\begin{itemize}
  \item FioName (\$B9B0+): device identifier
\end{itemize}

\textbf{Result:} FioStatus $=$ \$02 on success.

\medskip

\cmdentry{25}{Unmount}{Unmount disk}

Unmounts the active storage device.

\textbf{Parameters:} None.

\textbf{Result:} FioStatus $=$ \$02 on success.

\medskip

\cmdentry{26}{Pwd}{Print working directory}

Writes the current working directory path into the FioName buffer.

\textbf{Parameters:} None.

\textbf{Result:} FioName populated with the current path.


% ═══════════════════════════════════════════════════════════════════════════════
\section{File Metadata Buffer (\$BAB0--\$BB1F)}
\label{sec:fio-metadata}
% ═══════════════════════════════════════════════════════════════════════════════

The 112-byte metadata buffer at \$BAB0--\$BB1F is populated by filtered
directory reads.  It provides rich metadata for SID files, MIDI files, and
binary executables --- information that is not available from the filename
alone.

\memrange{BAB0}{BB1F}{File Metadata Buffer}{R}{0}

\begin{center}
\small
\begin{longtable}{@{}lllp{7.5cm}@{}}
\toprule
\textbf{Address} & \textbf{Name} & \textbf{Size} & \textbf{Description} \\
\midrule
\endfirsthead
\toprule
\textbf{Address} & \textbf{Name} & \textbf{Size} & \textbf{Description} \\
\midrule
\endhead
\$BAB0 & MetaType      & 1  & File type (same encoding as FioDirType) \\
\$BAB1--\$BAB2 & MetaSize & 2 & File size in bytes (16-bit LE) \\
\$BAB3--\$BAD2 & MetaTitle & 32 & Title (null-padded ASCII) \\
\$BAD3--\$BAF2 & MetaAuthor & 32 & Author name (null-padded ASCII) \\
\$BAF3--\$BB12 & MetaCopyright & 32 & Copyright notice (null-padded ASCII) \\
\$BB13--\$BB14 & MetaLoadAddr & 2 & Load address (BIN/SID, 16-bit LE) \\
\$BB15--\$BB16 & MetaInitAddr & 2 & Init address (SID only, 16-bit LE) \\
\$BB17--\$BB18 & MetaPlayAddr & 2 & Play address (SID only, 16-bit LE) \\
\$BB19 & MetaSongs     & 1  & Song count (SID subtunes / MIDI tracks) \\
\$BB1A--\$BB1B & MetaDuration & 2 & Duration in seconds (MIDI only, 16-bit LE) \\
\$BB1C & MetaGfxSpace  & 1  & GFX space type \\
\bottomrule
\end{longtable}
\end{center}

\begin{tipbox}
Read the metadata buffer immediately after a DirRead command.  The buffer
contents are overwritten by each subsequent DirRead.
\end{tipbox}


% ═══════════════════════════════════════════════════════════════════════════════
\section{Status and Error Code Summary}
\label{sec:fio-status}
% ═══════════════════════════════════════════════════════════════════════════════

\subsection{Status Codes (FioStatus, \$B9A1)}

\begin{center}
\small
\begin{tabular}{@{}clp{8cm}@{}}
\toprule
\textbf{Value} & \textbf{Name} & \textbf{Description} \\
\midrule
\$00 & Idle  & No command has been issued, or the controller has been reset \\
\$02 & OK    & The last command completed successfully \\
\$03 & Error & The last command failed; check FioErrCode for the reason \\
\bottomrule
\end{tabular}
\end{center}

\subsection{Error Codes (FioErrCode, \$B9A2)}

\begin{center}
\small
\begin{tabular}{@{}clp{8cm}@{}}
\toprule
\textbf{Value} & \textbf{Name} & \textbf{Description} \\
\midrule
0 & None       & No error \\
1 & Not Found  & The specified file or directory does not exist \\
2 & I/O Error  & A read or write operation failed at the host level \\
3 & End of Dir & All directory entries have been enumerated \\
4 & Disk Full  & Insufficient space on the storage device \\
5 & Not Mounted & The target storage device is not mounted \\
\bottomrule
\end{tabular}
\end{center}


% ═══════════════════════════════════════════════════════════════════════════════
\section{BASIC Examples}
\label{sec:fio-examples}
% ═══════════════════════════════════════════════════════════════════════════════

\subsection{Saving and Loading a Program}

\begin{lstlisting}
10 REM === SAVE a BASIC program ===
20 REM set filename "MYPROG"
30 F$ = "MYPROG" + CHR$(0)
40 FOR I=0 TO LEN(F$)-1
50   POKE $B9B0+I, ASC(MID$(F$,I+1,1))
60 NEXT I
70 REM set source range ($0300-$2000)
80 DOKE $B9A4, $0300 : REM FioSrc
90 DOKE $B9A6, $2000 : REM FioEnd
100 POKE $B9A0, $01  : REM execute Save
110 IF PEEK($B9A1)=$02 THEN PRINT "SAVED OK"
120 IF PEEK($B9A1)=$03 THEN PRINT "ERROR:"; PEEK($B9A2)
\end{lstlisting}

\begin{lstlisting}
10 REM === LOAD a BASIC program ===
20 F$ = "MYPROG" + CHR$(0)
30 FOR I=0 TO LEN(F$)-1
40   POKE $B9B0+I, ASC(MID$(F$,I+1,1))
50 NEXT I
60 DOKE $B9A4, $0300 : REM FioSrc (load destination)
70 POKE $B9A0, $02   : REM execute Load
80 IF PEEK($B9A1)=$02 THEN PRINT "LOADED"; DEEK($B9A8); "BYTES"
90 IF PEEK($B9A1)=$03 THEN PRINT "ERROR:"; PEEK($B9A2)
\end{lstlisting}

\subsection{Directory Listing}

\begin{lstlisting}
10 REM list all files in the save directory
20 POKE $B9A0, $03 : REM DirOpen
30 REM loop through entries
40 POKE $B9A0, $04 : REM DirRead
50 IF PEEK($B9A1)=$03 THEN PRINT "---END---" : END
60 REM read filename from buffer
70 N$ = ""
80 FOR I=0 TO 63
90   C = PEEK($B9B0+I)
100  IF C=0 THEN I=63 : GOTO 120
110  N$ = N$ + CHR$(C)
120 NEXT I
130 T = PEEK($B9AF) : REM file type
140 T$ = "???"
150 IF T=0 THEN T$ = "BAS"
160 IF T=1 THEN T$ = "SID"
170 IF T=2 THEN T$ = "BIN"
180 IF T=3 THEN T$ = "MID"
190 IF T=5 THEN T$ = "DIR"
200 PRINT T$; " "; N$
210 GOTO 40
\end{lstlisting}

\subsection{Playing a SID File}

\begin{lstlisting}
10 REM play a SID file named "COMMANDO"
20 F$ = "COMMANDO" + CHR$(0)
30 FOR I=0 TO LEN(F$)-1
40   POKE $B9B0+I, ASC(MID$(F$,I+1,1))
50 NEXT I
60 POKE $B9A0, $08 : REM SidPlay
70 IF PEEK($B9A1)=$02 THEN PRINT "PLAYING..."
80 REM press any key to stop
90 GET K$ : IF K$="" THEN 90
100 POKE $B9A0, $09 : REM SidStop
110 PRINT "STOPPED"
\end{lstlisting}

\subsection{Loading and Playing a MIDI File}

\begin{lstlisting}
10 REM load soundfont first
20 F$ = "GeneralUser.sf2" + CHR$(0)
30 FOR I=0 TO LEN(F$)-1
40   POKE $B9B0+I, ASC(MID$(F$,I+1,1))
50 NEXT I
60 POKE $B9A0, $15   : REM SfLoad
70 WAIT $A183, 1     : REM wait for WTS ready
80 PRINT PEEK($A184); " INSTRUMENTS LOADED"
90 REM now play a MIDI file
100 F$ = "BACH_FUGUE" + CHR$(0)
110 FOR I=0 TO LEN(F$)-1
120   POKE $B9B0+I, ASC(MID$(F$,I+1,1))
130 NEXT I
140 POKE $B9A0, $13  : REM MidPlay
150 IF PEEK($B9A1)=$02 THEN PRINT "MIDI PLAYING"
160 REM press any key to stop
170 GET K$ : IF K$="" THEN 170
180 POKE $B9A0, $14  : REM MidStop
190 PRINT "STOPPED"
\end{lstlisting}

\begin{notebox}
The BASIC interpreter provides higher-level commands for many of these
operations: \texttt{SAVE}, \texttt{LOAD}, \texttt{DIR}, \texttt{DEL},
\texttt{SIDPLAY}, \texttt{SIDSTOP}, \texttt{MUSIC}, and others.  Direct FIO
register access is most useful for operations not exposed by BASIC, for
programs that need fine-grained error handling, or for accessing the metadata
buffer.
\end{notebox}
